\begin{scanpage}{0035}{18}
\sectiontitle{3}{The Physical Optics Approach}
In the physical optics approach, the fringing field is sought as the electro-magnetic field created by the uniform part of the surface current
\[\mathbf j^0=\frac c{2\pi}[\mathbf n\mathbf H_0].\tag{3.01}\]
Let us recall that here $\mathbf n$ designates the external normal to the body's surface, and $\mathbf H_0$ designates the magnetic vector of the incident wave. First let us investigate the case $0<\phi_0<\alpha-\pi$, when the incident plane wave irradiates only one face of the wedge (Figure 5).

From Equation (3.01) it follows that the density of the uniform part of the current being excited on the irradiated face by plane waves (1.23) and (1.24) has the following components, respectively
\[j_z^0=\frac c{2\pi}E_{0z}\sin\phi_0e^{-ikx\cos\phi_0},\quad j_x^0=j_y^0=0\tag{3.02}\]
and
\[j_x^0=\frac c{2\pi}H_{0z}e^{-ikx\cos\phi_0},\quad j_y^0=j_z^0=0.\tag{3.03}\]
For the purpose of calculating the field radiated by this current, we will use the following integral representation of the Hankel function (see [16], p.\ 866)
\[H_0^{(1)}(\rho)=\frac1\pi\int_{-\delta+i\infty}^{\delta-i\infty}e^{i\rho\cos\beta}\,d\beta=\frac1{i\pi}\int_{-\infty}^{\infty}e^{i\rho\operatorname{ch}t}\,dt\quad(0\leq\delta\leq\pi).\tag{3.04}\]
Assuming here $\rho=kd$ and changing to a new integration variable $\zeta=d\operatorname{sh}t$, we obtain
\end{scanpage}
